\documentclass{article}
\usepackage{amsmath,amssymb}
\usepackage{xurl}
\usepackage[hidelinks]{hyperref}
\setlength{\emergencystretch}{2em}
% Shared commands for notes in docs/paper-gaps/.

\DeclareUnicodeCharacter{2080}{\ifmmode{}_{0}\else\textsubscript{0}\fi}
\DeclareUnicodeCharacter{2081}{\ifmmode{}_{1}\else\textsubscript{1}\fi}
\DeclareUnicodeCharacter{2082}{\ifmmode{}_{2}\else\textsubscript{2}\fi}
\DeclareUnicodeCharacter{2099}{\ifmmode{}_{n}\else\textsubscript{n}\fi}

\newcommand{\Lean}{\textsc{Lean}}
\ExplSyntaxOn
\NewDocumentCommand{\leanid}{m}{
  \ifmmode
    \textsf{\detokenize{#1}}
  \else
    \group_begin:
    \urlstyle{sf}
    \tl_set:Nn \l_tmpa_tl {#1}
    \tl_replace_all:Nnn \l_tmpa_tl {\_} {_}
    \exp_args:NV \path \l_tmpa_tl
    \group_end:
  \fi
}
\ExplSyntaxOff
\providecommand{\tr}{\operatorname{tr}}
\newcommand{\tnleanIssueBase}{https://github.com/LionSR/TNLean/issues/}
\newcommand{\tnleanPrBase}{https://github.com/LionSR/TNLean/pull/}
\newcommand{\tnleanDocBase}{https://sirui-lu.com/TNLean/docs/}
\newcommand{\ghissue}[1]{\href{\tnleanIssueBase#1}{GitHub issue \##1}}
\newcommand{\ghpr}[1]{\href{\tnleanPrBase#1}{PR \##1}}
\newcommand{\doclink}[1]{\href{\tnleanDocBase#1.html}{\path{#1}}}

% Dirac notation and the identity operator, as in blueprint/src/macros/common.tex.
% \providecommand keeps note-local definitions authoritative.
\usepackage{dsfont}
\providecommand{\ket}[1]{|#1\rangle}
\providecommand{\bra}[1]{\langle #1|}
\providecommand{\braket}[2]{\langle #1 | #2 \rangle}
\providecommand{\ketbra}[2]{|#1\rangle\!\langle #2|}
\providecommand{\Id}{\mathds{1}}
\providecommand{\one}{\mathds{1}}
\providecommand{\C}{\mathbb{C}}
\providecommand{\MN}[1]{M_{#1}(\C)}
\providecommand{\MD}{M_D(\C)}

% Verdict marker: \gapnote{<kind>}{<status>}. Kind is the policy
% classification (clarification, local-correction, scope-restriction,
% unfaithful, false-source, open-gap); status is open, wip (elimination
% underway), resolved, or historical. Machine-read by the site index;
% prints a verdict line.
\newcommand{\gapnote}[2]{\par\noindent\textbf{Verdict:} #1 (#2).\par}


\title{A Gapless Short-Range Parent Hamiltonian of a Canonical MPS}
\author{}
\date{2026-10-02}

\begin{document}
\maketitle
\gapnote{false-source}{resolved}

\section{The unrestricted source assertion}

The parent-interaction definition in
\cite[Section~IV.C]{Cirac2021Matrix} permits any positive operator on $L$
sites whose kernel is the local matrix product space $G_L$. The subsequent
gap theorem states that all MPS parent Hamiltonians have a positive gap
uniform in the chain length. There is no interaction-range qualification
in that theorem. The local source locations are
\path{Papers/2011.12127/TN-Review-main.tex}:1996--1999 and 2183--2187.
Comparable uniqueness and ground-space statements explicitly specify their
interaction lengths, so a sufficient-range hypothesis cannot be inserted
silently into this assertion.

The example below has a nonzero range-two parent interaction with exactly
the kernel required by the source definition. Its tensor is in canonical
form, with three inequivalent normalized normal blocks of positive bond
dimension and nonzero weights. Nevertheless its positive excitation
energies approach zero as the ring length tends to infinity.
The counterexample is verified by the declarations listed below: the local
parent kernel, the normalized normal-block data, the periodic Fourier
eigenvector, and failure of an eventual uniform gap are all proved.

\section{Three normal product-state blocks}

Let $\alpha=1/\sqrt2$, and define the two physical matrices
\begin{align}
    B^0&=\operatorname{diag}(1,0,\alpha),\notag\\
    B^1&=\operatorname{diag}(0,1,\alpha).
    \label{eq:short_range_parent_tensor}
\end{align}
Thus $B=A^0\oplus A^1\oplus A^+$ has three bond-one blocks
$A^0=(1,0)$, $A^1=(0,1)$, and $A^+=(\alpha,\alpha)$, each with weight one.
Their product vectors are $|0\rangle^{\otimes N}$,
$|1\rangle^{\otimes N}$, and $|+\rangle^{\otimes N}$, where
$|+\rangle=(|0\rangle+|1\rangle)/\sqrt2$.

For each block, the sum of squared moduli of its two letters is one.
Its transfer map on $\mathbb C$ is therefore the identity, which has
spectral radius one and is primitive on the cone of nonnegative scalars.
The block is normal and normalized. Similarity is trivial in bond
dimension one, and no pair of these two-component vectors differs by a
phase. All weights are nonzero. Thus
(\ref{eq:short_range_parent_tensor}) is a canonical tensor of bond
dimension three, with no zero or empty summand.

The simultaneous word tuples at length two include
\begin{align}
    w_{00}&=(1,0,1/2),\notag\\
    w_{01}&=(0,0,1/2),\notag\\
    w_{11}&=(0,1,1/2).
    \label{eq:short_range_parent_word_tuples}
\end{align}
They span $\mathbb C^3$, since the coordinate vectors are
$e_0=w_{00}-w_{01}$, $e_1=w_{11}-w_{01}$, and $e_2=2w_{01}$.
At length one there are only two word tuples, so the simultaneous span
cannot have dimension three. The simultaneous injectivity length is
therefore $S=2$.

\section{The exact two-site parent interaction}

For an arbitrary boundary matrix $X\in M_3(\mathbb C)$, the diagonal
form in (\ref{eq:short_range_parent_tensor}) gives
\begin{align}
    \Gamma_2^B(X)
      &=X_{00}|00\rangle+X_{11}|11\rangle+X_{22}|++\rangle.
    \label{eq:short_range_parent_boundary_map}
\end{align}
It follows that
\begin{align}
    G_2(B)
      &=\operatorname{span}\{|00\rangle,|11\rangle,|++\rangle\}
       =\operatorname{Sym}^2(\mathbb C^2).
    \label{eq:short_range_parent_local_space}
\end{align}
Indeed, $2|++\rangle-|00\rangle-|11\rangle=|01\rangle+|10\rangle$.
The orthogonal complement is the singlet line generated by
$|\psi^-\rangle=(|01\rangle-|10\rangle)/\sqrt2$.

Let $F$ exchange the two physical sites. The canonical parent interaction
is the nonzero positive projection
\begin{align}
    h_2&=|\psi^-\rangle\langle\psi^-|=\frac{I-F}{2}
       =\frac12
       \begin{pmatrix}
         0&0&0&0\\
         0&1&-1&0\\
         0&-1&1&0\\
         0&0&0&0
       \end{pmatrix}.
    \label{eq:short_range_parent_singlet}
\end{align}
Its kernel is exactly (\ref{eq:short_range_parent_local_space}). In terms
of spin-$1/2$ operators, $h_2=\frac14I-\mathbf S_1\cdot\mathbf S_2$.
Consequently its periodic sum is the isotropic ferromagnetic Heisenberg
chain with ground energy zero.

For comparison with the primary literature, the isotropic case of
Koma--Nachtergaele~\cite{Koma1997Spectral}, equations (2.1)--(2.3), has this local
singlet projection. Their Lemma~4, equation (3.13), gives the same
interior one-magnon recurrence. Their finite-volume boundary convention
is open; the periodic formula below is derived directly.

\section{Positive one-magnon eigenvalues on the ring}

Fix $N\geq3$ and index the sites by $\mathbb Z/N\mathbb Z$. Define
$H_N=\sum_j(h_2)_{j,j+1}$. Let $|j\rangle$ denote the computational
basis vector with a single spin in state $1$ at site $j$ and every other
spin in state $0$. These $N$ vectors are orthonormal.

Every term not adjacent to $j$ annihilates $|j\rangle$. The two adjacent
terms in (\ref{eq:short_range_parent_singlet}) give
\begin{align}
    H_N|j\rangle
       &=|j\rangle-\frac12|j-1\rangle-\frac12|j+1\rangle.
    \label{eq:short_range_parent_magnon_action}
\end{align}
Put $\theta_N=2\pi/N$ and
$v_N=\sum_{j=0}^{N-1}e^{i\theta_Nj}|j\rangle$. Its squared norm is
$N$, so it is nonzero. The periodic boundary condition is compatible with
these coefficients because $e^{i\theta_NN}=1$. Reindexing the two shifted
sums in (\ref{eq:short_range_parent_magnon_action}) yields
\begin{align}
    H_Nv_N
       &=\left(1-\frac{e^{i\theta_N}+e^{-i\theta_N}}2\right)v_N
        =\varepsilon_N v_N,\notag\\
    \varepsilon_N&=1-\cos(2\pi/N).
    \label{eq:short_range_parent_magnon_eigenvalue}
\end{align}
Since $0<2\pi/N<\pi$, the eigenvalue is strictly positive. For
$w\in\ker H_N$, self-adjointness gives
$\varepsilon_N\langle w,v_N\rangle=\langle H_Nw,v_N\rangle=0$.
Thus $v_N\perp\ker H_N$, without any assumption about the dimension
of the ground space.

As $N\to\infty$, $2\pi/N\to0$ and therefore $\varepsilon_N\to0$.
If a positive constant $\gamma$ satisfied the source's uniform gap
inequality on every ring, (\ref{eq:short_range_parent_magnon_eigenvalue})
would imply
$\gamma\|v_N\|\leq\|H_Nv_N\|=\varepsilon_N\|v_N\|$.
Since $v_N\neq0$, this forces $\gamma\leq\varepsilon_N$ for every
$N\geq3$, a contradiction.

\section{Formal verification}

The normal-block data and exact simultaneous injectivity length two are
formally proved.\footnote{Formal declarations:
    \leanid{MPSTensor.isCPSVBasisOfNormalTensors_shortRangeHeisenbergTensor},
    \leanid{MPSTensor.wordTupleSpanTop_shortRangeHeisenbergBlock_two}, and
    \leanid{MPSTensor.not_wordTupleSpanTop_shortRangeHeisenbergBlock_one}
    in \doclink{TNLean/MPS/ParentHamiltonian/ShortRangeHeisenbergTensor}.}
The two-site ground-space criterion is equality of the $01$ and $10$
coefficients, and the parent interaction is the singlet projection.\footnote{Formal declarations:
    \leanid{MPSTensor.mem_groundSpaceES_shortRangeHeisenbergTensor_two_iff} and
    \leanid{MPSTensor.parentInteractionES_shortRangeHeisenbergTensor_two_apply}
    in the same module.}

The Fourier eigenvector, its positive eigenvalue limit, and failure of an
eventual uniform gap are formally proved.\footnote{Formal declarations:
    \leanid{SpinChain.oneMagnon_edge_sum},
    \leanid{ZMod.cycleLaplacian_stdAddChar},
    \leanid{ZMod.cycleFourierGap_pos},
    \leanid{ZMod.tendsto_cycleFourierGap_zero},
    \leanid{MPSTensor.parentHamiltonianES_shortRangeHeisenbergMagnon}, and
    \leanid{MPSTensor.not_exists_eventual_uniform_gap_shortRangeHeisenbergTensor_two}.
    The final two declarations are in
    \doclink{TNLean/MPS/ParentHamiltonian/ShortRangeHeisenbergGap}.}

\section{The range-qualified theorem remains valid}

This example has interaction range $R=2$ and simultaneous injectivity
length $S=2$. It does not satisfy the sufficient condition $R\geq S+1$.
It therefore leaves intact the proved gap theorem at every range
$R\geq S+1$, as well as the automatic sufficient dimension bound
$R\geq3D^5$ for canonical tensors. The corresponding prescribed-range
comparison is documented in
\path{docs/paper-gaps/cpgsv21_block_parent_interaction_range.tex}.
That note resolves the passage from a chosen large range to the specified
PGVWC07 and simultaneous-injectivity ranges; it does not prove the
unrestricted assertion refuted here.

The earlier construction in
Fern\'andez-Gonz\'alez et al.~\cite{FernandezGonzalez2015Gapless},
Section~2.3, Definition~3 and Remark~1, explicitly
uses a jointly injective blocked tensor and specifies a sufficient
interaction length. Its gap theorem does not supply the broader claim
for an arbitrary shorter range. The present example uses the exact
source-defined parent kernel throughout, rather than the kernel of an
uncle Hamiltonian.

\bibliographystyle{alpha}
\bibliography{references}
\end{document}
